\section{Problema y pregunta de investigación}
El proyecto estudiará cómo mejorar Qwen3-1.7B en problemas cuantitativos con
unidades mediante postentrenamiento. Partimos del entorno construido en el
desafío, que genera tareas en español y verifica sus resultados automáticamente,
para comparar ajuste supervisado (SFT), aprendizaje por refuerzo con recompensa
verificable (RL) y optimización directa de preferencias (DPO). Esta entrega
formula la investigación y su protocolo tentativo, sin presentar entrenamientos
del proyecto como ejecutados.

La pregunta principal es: \emph{¿cómo influye realizar SFT antes de RL en la
eficiencia de aprendizaje, frente a solo RL, solo SFT y DPO, bajo un presupuesto
de cómputo comparable?} SFT podría facilitar que el modelo produzca soluciones
válidas para recibir recompensa, aunque también podría consumir presupuesto
sin mejorar la generalización. Por eso distinguiremos adaptación al formato,
exactitud numérica y calidad del razonamiento.

La hipótesis de trabajo es que SFT previo mejora el rendimiento de validación
por unidad de cómputo frente a iniciar directamente con RL. Se contrastará
también la alternativa de que la mejora se concentre en el formato o desaparezca
al contabilizar el costo de SFT. La revisión siguiente fundamenta esa comparación
y delimita lo que sus antecedentes permiten esperar.

\section{Antecedentes y relación con la propuesta}
Stojanovski et al. (2025) presentan Reasoning Gym~\cite{reasoninggym}, una
biblioteca de tareas generadas por código con dificultad controlable y
verificación automática. Su diseño respalda utilizar un entorno propio para
producir señal de recompensa y estudiar aprendizaje por dificultad, pero sus
resultados no demuestran que nuestro conjunto de unidades tenga volumen o
calibración suficientes para entrenar.

GSM-Symbolic, de Mirzadeh et al. (2024)~\cite{gsmsymbolic}, muestra variación
del desempeño al cambiar valores y condiciones de problemas matemáticos.
Esto motiva conservar variantes relacionadas en la misma partición y evaluar
por familia y nivel, evitando atribuir generalización a un promedio global.
Nuestro experimento compartirá fórmulas entre particiones, por lo que estudiará
nuevas magnitudes y representaciones, no nuevos dominios de razonamiento.

Rafailov et al. (2023)~\cite{dpo} formulan DPO como entrenamiento directo con
pares preferidos y rechazados, sin ajustar un modelo explícito de recompensa
ni muestrear nuevas respuestas durante la actualización. Sus experimentos de
preferencias motivan esta alternativa a RL, pero no establecen superioridad en
tareas con verificación numérica. Aquí las preferencias procederán de soluciones
verificadas y revisadas, y el reducido número inicial de pares será una limitación
a controlar.

DeepSeekMath, de Shao et al. (2024)~\cite{deepseekmath}, introduce GRPO, que
compara las recompensas de varias respuestas al mismo problema y prescinde del
modelo de valor adicional de PPO. Sus resultados de RL parten de un modelo
previamente ajustado con SFT, de modo que no resuelven directamente nuestra
comparación con RL desde el checkpoint inicial. Proponemos GRPO para ambas
ramas de RL, sujeto a comprobar memoria y tiempo en el hardware disponible.
La menor carga respecto de PPO no garantiza que quepa en una GPU de 6 GiB.

Finalmente, Aleithan et al. (2024)~\cite{swebenchplus} muestran mediante revisión
manual que superar pruebas automáticas de programación puede ocultar soluciones
problemáticas. Tomamos esa advertencia metodológica, sin trasladar sus tasas a
nuestro dominio: acompañaremos la recompensa numérica con una auditoría humana
de respuestas exitosas. Así, el estudio relacionará desempeño medido con la
validez de las soluciones que el modelo aprende a producir.

\section{Entorno y datos disponibles}
El entorno del desafío~\cite{challenge} contiene tres familias, entendidas como
tipos de problema definidos por una fórmula: velocidad $v=d/t$, energía cinética
$E=mv^2/2$ y luminosidad $L=4\pi d^2F$. Se proporcionan fórmula y equivalencias,
por lo que el modelo debe componer operaciones y conversiones sin recordar
constantes ni utilizar herramientas. Cada problema admite datos directos o
sumas de distancias, masas o flujos y se presenta en tres niveles: sin
conversiones, con una conversión de entrada o con dos entradas convertidas y,
en velocidad y energía difíciles, también la salida.

Los 60 problemas físicos generan $60\times3=180$ tareas. Por cada familia se
asignan 10 problemas a entrenamiento, cuatro a validación y seis a prueba,
manteniendo juntas sus variantes. Esto produce 90, 36 y 54 tareas respectivamente.
Por ejemplo, una tarea proporciona distancias de 4,08 y 1,02 km y un tiempo
de 88 s, y pide $(d_1+d_2)/t$ en m/s. La salida \verb|57.95 m/s| es aceptable
frente al resultado esperado $5100/88$ m/s.

La función \texttt{reward(prompt, completion)} devuelve 1 si extrae un único
resultado finito, de magnitud compatible y dentro del 1\% de la referencia,
tras convertirlo al Sistema Internacional, y 0 en caso contrario. Las referencias
se calculan por código y no se envían al modelo. El extractor puede rechazar
formatos correctos y actualmente admite unidades equivalentes distintas de la
solicitada, por lo que su recompensa no certifica todos los pasos ni obliga
siempre a realizar la conversión de salida.

Disponemos de 121 soluciones SFT generadas por Qwen3-8B y 35 pares DPO,
todos de entrenamiento y con procedencia registrada. Su revisión fue realizada
por agentes, no certificada por humanos, y energía y luminosidad difíciles
conservan solo una solución SFT cada una. Estos datos servirán como punto de
partida, sujetos a revalidación y ampliación antes de fijar los experimentos.
La línea base previa de Qwen3-1.7B fue 51,85\% en pass@1 y 68,52\% en pass@3,
con cobertura de extracción de 95,68\%. Se obtuvo con Ollama y pesos GGUF
cuantizados, y deberá repetirse si cambia la configuración para entrenar.

\section{Comparación experimental propuesta}
Se plantea un estudio de métodos existentes con implementaciones públicas.
Todas las rutas partirán del mismo checkpoint inicial de Qwen3-1.7B, entendido
como el modelo antes de nuestro ajuste adicional, y usarán los mismos enunciados,
particiones y evaluación. La tabla resume qué información recibe cada estrategia
y evita confundir la comparación de modelos del desafío con esta comparación
de entrenamientos de un único modelo.

\begin{center}
\begin{tabular}{lp{.68\linewidth}}
\toprule
Ruta & Procedimiento \\
\midrule
Control & Evaluar el checkpoint inicial sin ajuste adicional. \\
Solo SFT & Ajustar con las soluciones correctas de entrenamiento. \\
SFT $\rightarrow$ RL & Aplicar SFT y continuar con GRPO y recompensa del verificador. \\
Solo RL & Aplicar GRPO desde el checkpoint inicial, sin SFT propio previo. \\
DPO & Ajustar desde el checkpoint inicial con pares preferidos/rechazados y referencia inicial fija. \\
\bottomrule
\end{tabular}
\end{center}

Primero se medirá un piloto de memoria y tiempo, a partir del cual se fijará un
presupuesto $B$ de entrenamiento por ruta en hardware común. SFT seguido de RL
contabilizará ambas fases dentro de $B$, mientras el control no consume ese
presupuesto. Se registrarán horas de GPU, tokens y evaluaciones por separado,
identificando también el costo compartido de generar datos. Así, una ventaja
de SFT previo no podrá atribuirse simplemente a disponer de más entrenamiento.

La sensibilidad principal será la fracción de $B$ dedicada al SFT previo:
se propone comparar 25\% y 50\%, dejando el resto para RL, si el piloto permite
ejecutar ambas condiciones. Los demás hiperparámetros de RL se mantendrán
comunes entre las ramas para estudiar el efecto de su inicialización. Los
checkpoints se seleccionarán con validación, nunca con el conjunto de prueba.
Si los datos requieren ampliación, se hará antes de fijar la comparación y
se conservará la separación por problema físico.

\section{Evaluación, recursos y próximos hitos}
Se medirán pass@1 y pass@3 con tres respuestas por tarea, por familia y nivel.
Pass@1 estima la proporción de respuestas correctas y pass@3 la proporción de
tareas resueltas al menos una vez entre tres intentos. Se conservarán cobertura,
errores numéricos y respuestas no interpretadas para distinguir capacidad de
cálculo de cumplimiento del formato. El muestreo será común a las cinco rutas,
con semillas 42--44 y límites de generación fijados con validación.

La hipótesis se contrastará con curvas de pass@1 en validación frente al cómputo
acumulado, incluyendo el SFT previo, y con el rendimiento final a igual $B$.
Como complemento se auditará manualmente una muestra de hasta tres éxitos por
familia, nivel y ruta, elegidos aleatoriamente cuando existan. Se revisarán datos,
conversiones, operaciones y conclusión, informando por estrato el denominador y
la proporción de éxitos genuinos sin corregir silenciosamente la recompensa automática.

Antes de entrenar se precisarán y probarán las políticas de formatos, unidades
y respuestas contradictorias. Se revalidarán SFT y DPO si cambia la recompensa
y se reservará un test independiente del ajuste, porque el test del desafío ya
fue observado al diagnosticar límites. Todos los resultados finales se evaluarán
con la misma versión final del verificador, preservando el historial de problemas,
correcciones y reentrenamientos. La dependencia entre variantes y la escasez de
algunas celdas limitarán las conclusiones estadísticas.

El recurso inicial confirmado es una RTX 3050 Laptop de 6 GiB. El piloto probará
ajuste de pocos parámetros mediante adaptadores LoRA y respuestas acotadas,
sin asumir que la inferencia previa demuestra viabilidad de RL. Otro colaborador
podría aportar cómputo, pero no se considera disponible todavía. Si se cambia de
GPU se repetirá el piloto y se compararán las rutas en condiciones equivalentes,
sin tratar horas de equipos distintos como recursos idénticos. Las semillas de
entrenamiento adicionales dependerán del cómputo que quede disponible.

Para el avance del 25 de noviembre se prevén el piloto, la adaptación del entorno,
la nueva línea base y primeras curvas de entrenamiento. Para el cierre del
9 de diciembre se completarán las cinco rutas, sensibilidad, auditoría y análisis,
junto con paper, poster y video. Esta propuesta delimita esos experimentos y
sus riesgos sin adelantar mejoras que todavía no se han medido.

\clearpage
\begingroup\raggedright
\begin{thebibliography}{9}
\bibitem{reasoninggym} Z. Stojanovski, O. Stanley, J. Sharratt, R. Jones,
A. Adefioye, J. Kaddour y A. K\"opf. Reasoning Gym: Reasoning Environments
for Reinforcement Learning with Verifiable Rewards. 2025.
\url{https://arxiv.org/abs/2505.24760}.
\bibitem{gsmsymbolic} I. Mirzadeh, K. Alizadeh, H. Shahrokhi, O. Tuzel,
S. Bengio y M. Farajtabar. GSM-Symbolic: Understanding the Limitations of
Mathematical Reasoning in Large Language Models. 2024.
\url{https://arxiv.org/abs/2410.05229}.
\bibitem{dpo} R. Rafailov, A. Sharma, E. Mitchell, S. Ermon,
C. D. Manning y C. Finn. Direct Preference Optimization: Your Language Model
Is Secretly a Reward Model. 2023. \url{https://arxiv.org/abs/2305.18290}.
\bibitem{deepseekmath} Z. Shao, P. Wang, Q. Zhu, R. Xu, J. Song, X. Bi,
H. Zhang, M. Zhang, Y. K. Li, Y. Wu y D. Guo. DeepSeekMath: Pushing the Limits
of Mathematical Reasoning in Open Language Models. 2024.
\url{https://arxiv.org/abs/2402.03300}.
\bibitem{swebenchplus} R. Aleithan, H. Xue, M. M. Mohajer, E. Nnorom,
G. Uddin y S. Wang. SWE-Bench+: Enhanced Coding Benchmark for LLMs. 2024.
\url{https://arxiv.org/abs/2410.06992}.
\bibitem{challenge} F. Fajardo, S. Palma, B. Prieto y R. Ugarte.
Razonamiento cuantitativo con unidades: desafío de investigación. IIC-3694,
Pontificia Universidad Católica de Chile, 2026. Informe y datos disponibles en
\href{https://github.com/spmasabeu/research-challenge-ia-gen-202602/tree/e99f9e922b52bb56993d87e683f6a1bc217e7b36/deliveries/challenge-clean}{el repositorio del desafío}.
\end{thebibliography}
\endgroup
